\section{Simulation: bias, standard errors, and coverage}\label{sec:simulation}
The simulation asks whether an uncertainty calculation for the chosen
complete-label coefficient matches its repeated error. Each panel fixes
the item labels, loading matrices, controls, and any observed outcome.
Across repetitions, item predictions are generated by independent flips
with probability $0.15$, and each item is independently verified with
probability $0.70$. The audit probabilities and predictions are known when
applying equations~\eqref{eq:polynomialaudit}--\eqref{eq:auditratiose}.
This experiment isolates label prediction and verification; it does not
sample a new population of records in each repetition.

There are $200$ or $800$ records. The diffuse design uses $200$ items;
record-specific multinomial counts depend on an observed characteristic
and item features. The concentrated design gives one item three times the
baseline popularity of all other items combined. A third design uses only
$30$ items. Counts are converted to visit shares; a second aggregation
reweights them by item-specific durations. We examine an aggregate outcome,
an aggregate covariate with an observed outcome, and aggregates on both
regression sides. Appendix~\ref{app:simulation} gives the construction.
The complete-label coefficient is calculated separately for every fixed
panel and regression role.

The record HC method attaches HC0 SEs to the uncorrected proxy regression.
Those intervals concern the proxy fit under its usual sampling
interpretation; here we evaluate their use for the complete-label target.
The item-audit method corrects the regression moments and estimates their
joint covariance. An oracle version uses the same corrected estimate and
the true conditional moment covariance, with the regression derivative
evaluated at the true moments. It isolates variance estimation from
linearization and distributional failures and is not a feasible comparator
for an application without complete labels.

\begin{table}[htbp]
\centering\small
\caption{Regression inference with reused labels in simulated panels of 800 records}
\label{tab:regressioninference}
\begin{tabular}{lllrrrr}
\toprule
Role & Loadings & Method & Bias & SD & Mean SE & Cover (\%) \\
\midrule
Outcome & Diffuse & Record HC & -0.101 & 0.020 & 0.003 & 0.0 \\
Outcome & Diffuse & Item audit & 0.000 & 0.015 & 0.014 & 93.9 \\
Outcome & Concentrated & Record HC & -0.076 & 0.103 & 0.004 & 0.2 \\
Outcome & Concentrated & Item audit & 0.002 & 0.070 & 0.030 & 89.9 \\
Covariate & Diffuse & Record HC & 0.351 & 0.090 & 0.014 & 0.0 \\
Covariate & Diffuse & Item audit & 0.004 & 0.040 & 0.038 & 93.2 \\
Covariate & Concentrated & Record HC & -0.379 & 1.288 & 0.020 & 0.0 \\
Covariate & Concentrated & Item audit & -0.097 & 0.524 & 0.045 & 89.9 \\
Both & Diffuse & Record HC & -0.001 & 0.048 & 0.008 & 26.4 \\
Both & Diffuse & Item audit & 0.002 & 0.026 & 0.023 & 92.2 \\
Both & Concentrated & Record HC & 0.009 & 0.035 & 0.008 & 36.2 \\
Both & Concentrated & Item audit & 0.000 & 0.018 & 0.016 & 90.9 \\
\bottomrule
\end{tabular}

\par\smallskip
\begin{minipage}{\linewidth}\footnotesize
Each row summarizes \SimulationReplications{} repetitions of predictions and
independent item verification in one fixed panel. Outcome, covariate, and
both denote the roles of label-derived measures. There are 200 items; the
concentrated design gives one item a large common loading. Bias and empirical
SD are relative to the complete-label panel coefficient. Mean SE averages
the reported HC0 or item-audit delta SE. Coverage uses a normal 1.96 cutoff.
The largest binomial Monte Carlo SE across the simulation is
\SimulationMaxMCSE{} percentage points. No coefficient or SE is unavailable
in the displayed designs.
\end{minipage}
\end{table}
\FloatBarrier

In the diffuse $800$-record panels, item-audit coverage is
\SimOutcomeDiffuseCoverage{}\% for an aggregate outcome,
\SimCovariateDiffuseCoverage{}\% for an aggregate covariate, and
\SimBothDiffuseCoverage{}\% when both sides are constructed
(Table~\ref{tab:regressioninference}). In the covariate case, the
uncorrected coefficient's bias is \SimCovariateDiffuseNaiveBias{}, compared
with \SimCovariateDiffuseAuditBias{} after moment correction. Bias and
variance estimation thus answer distinct questions even in this favorable
loading pattern.

Concentration makes the delta approximation less reliable. In the
concentrated covariate example, the corrected coefficient has empirical SD
\SimCovariateConcentratedSD{}, while its mean estimated SE is
\SimCovariateConcentratedSE{} and coverage is
\SimCovariateConcentratedCoverage{}\%. Occasional small estimated
denominators produce large coefficient errors. Even the oracle delta
interval covers only \SimCovariateConcentratedOracleCoverage{}\% in this
design. With $30$ items and both sides constructed, feasible coverage is
\SimFewBothCoverage{}\%. These are results for the stated panels and error
process. They show why an exact, unbiased moment covariance is insufficient
to guarantee normal coefficient intervals. The smaller panels and oracle
comparisons appear in Appendix~\ref{app:simulation}.
